\documentclass[
	pdftex,				% PDFTex verwenden
	a4paper,			% A4 Papier
%	oneside,			% Einseitig
	twoside,			% Zweiseitig
%	bibtotocnumbered,	% Literaturverzeichnis nummeriert einfügen
%	idxtotoc,			% Index ins Verzeichnis einfügen
%	halfparskip,		% Europäischer Satz mit abstand zwischen Absätzen
%	chapterprefix,		% Kapitel anschreiben als Kapitel
	headsepline,		% Linie nach Kopfzeile
%	footsepline,		% Linie vor Fusszeile
	11pt,				% Größere Schrift, besser lesbar am Bildschrim
	openright,			% Anfang immer auf rechter Seite 
%	smallheadings,		% Kleinere Überschrifen
	headinclude,		% Header mit zum Satzspiegel rechnen
  ]{scrreprt}

% ----------------------------------------
% Sprach Erweiterungen einbinden
% ----------------------------------------

% Paket für Links und Verweise innerhalb des PDF Dokuments
\usepackage[
	pdftitle={},
	pdfauthor={},
	pdfcreator={},
	pdfsubject={},
	pdfkeywords={},
	plainpages=false,
	pdfpagelabels,
	bookmarksnumbered=true,
	bookmarksopen=true,
	pdfpagemode=UseOutlines
]{hyperref}

\usepackage{xurl}
% URLs im Dokument einbetten

% Paket für Deutsche Sprach Unterstützung
\usepackage[english, ngerman]{babel}
\selectlanguage{english}
 
% Latex Eingabedateien im UTF8 Format
\usepackage[utf8x]{inputenc}

% Pakete um Latin1 Zeichnensätze verwenden zu können und die dazu passenden Schriften
\usepackage[T1]{fontenc}

% Schriftart mit deutschen Umlauten
\usepackage{lmodern} 

% Sonderzeichen
\usepackage{textcomp}

% Pakete um Grafiken einbetten zu können
\usepackage[pdftex]{graphicx}
\usepackage{wrapfig}

% Zeilenabstand anpassen
% \onehalfspacing, \doublespacing, \singlespacing
\usepackage{setspace}

% Paket um Textteile drehen zu können
\usepackage{rotating}

% Für die Steuerung von Linien in Tabellen
% \toprule, \midrule, \bottomrule
\usepackage{booktabs}

% Erlaubt mehrseitige Tabellen
\usepackage{longtable}

% Erlaubt in einer Tabelle die Felder mehrerer Zeilen zusammenzufassen
\usepackage{multirow}

% Verbessert einige Fremdpakete und das Zusammenspiel von einigen Paketen mit KOMA-Script
\usepackage{scrhack}

% Einbinden von Seiten aus PDFs
\usepackage{pdfpages}

% \usepackage[backend=biber, style=authoryear, url=true]{biblatex}
\usepackage[
  backend=biber,
  style=authoryear,
  url=true,
  doi=true
]{biblatex}

\renewcommand{\mkbibparens}[1]{\bibopenbracket#1\bibclosebracket}

\DeclareDelimFormat{nameyeardelim}{\addspace}
\let\cite\parencite

\addbibresource{database.bib}
% ------------------------------------------
% Farben
% ------------------------------------------

% Package für Farbdefinitionen
\usepackage{xcolor}

% Link Farbe innerhalb des Dokuments
\definecolor{LinkColor}{rgb}{0,0,0.6}
% Code Box
\definecolor{CodeKeywords}{rgb}{0,0,0}
\definecolor{CodeComments}{rgb}{0,0.6,0}
\definecolor{CodeNumbers}{rgb}{0.5,0.5,0.5}
\definecolor{CodeStrings}{rgb}{0.58,0,0.82}
\definecolor{CodeBackground}{rgb}{1.0,1.0,1.0}

% Code & Theorem Hintergrund
\definecolor{BoxBackground}{rgb}{0.95,0.95,0.95}

% -------------------------------------------
% Dokument Einstellungen
% -------------------------------------------

% Benutzerdefinierte Größe des Untertitels (Bilder usw.)
\usepackage[size=normalsize]{caption}

% Zeilenabstand
\onehalfspacing
\setlength{\parskip}{6pt}
\setlength{\parindent}{0pt}

% Schusterjungen und Hurenkinder vermeiden
\clubpenalty = 10000	% Schusterjunge
\widowpenalty = 10000	% Hurenkind

% Zeilenumbruch bei Bildbeschreibungen
\setcapindent{1em}

% Link Farben innerhalb des Dokuments
\hypersetup{
	colorlinks=true,
	linkcolor=black,
	citecolor=black,
	filecolor=LinkColor,
	menucolor=LinkColor,
	urlcolor=LinkColor}

% Bis zu welcher Strukturtiefe soll nummeriert werden
\setcounter{secnumdepth}{2}

% Fussonotennummerierung durch das gesamte Dokument
\usepackage{chngcntr}
\counterwithout{footnote}{chapter}

% Spezielle Schrift verwenden
\renewcommand{\encodingdefault}{T1}
%\renewcommand{\familydefault}{pfr}
%\renewcommand{\sfdefault}{pfr}

% Spezielle Schrift im Koma-Script setzen
%\setkomafont{llabel}{\normalfont\bfseries}
%\setkomafont{sectioning}{\sffamily\bfseries}
%\setkomafont{captionlabel}{\sffamily\bfseries}
%\setkomafont{pagehead}{\sffamily\small\itshape}
%\setkomafont{descriptionlabel}{\sffamily\bfseries}

% Nur Punkte in Aufzählungen verwenden (bis Ebene 3)
\renewcommand*\labelitemii{\textbullet}
\renewcommand*\labelitemiii{\textbullet}

% -----------------------------------------
% Seiten Layout Einstellungen
% -----------------------------------------

% Paket zur einfacheren Manipulation von Kopf- und Fußzeile
\usepackage{fancyhdr}
\pagestyle{fancy}
%\pagestyle{empty}

% Seitenabmessungen
\setlength{\topmargin}{0pt}
\setlength{\voffset}{0cm}
\setlength{\textheight}{20.5 cm}
%\setlength{\headsep}{0.7cm}
\setlength{\evensidemargin}{1.15cm}
\setlength{\oddsidemargin}{0.2cm}
\setlength{\textwidth}{14.7cm}
\setlength{\headwidth}{14.7cm}
\setlength{\skip\footins}{15mm}
%\setlength{\footskip}{2cm}

% Kapitel Anzeige im Header
\renewcommand{\chaptermark}[1]{\markboth{\MakeUppercase{\thechapter\ #1}}{}}

% Kopf- und Fusszeilen formatieren
\fancypagestyle{normal}{
  \fancyhf{}
  \renewcommand\headrulewidth{0.4pt}
  \fancyhead[RE]{\leftmark}
  \fancyhead[LO]{\rightmark}
  \fancyhead[RO,LE]{\thepage}
}

% Anfang eines neuen Kapitels etc.
\fancypagestyle{plain}{
  \renewcommand\headrulewidth{0.0pt}
  \fancyhf{}
  \fancyfoot[CO]{\thepage}
}

% Header für das Inhaltsverzeichnis
\fancypagestyle{ihv}{
  \fancyhf{}
  \renewcommand\headrulewidth{0.4pt}
  \fancyhead[RE,LO]{CONTENTS}
  \fancyhead[RO,LE]{\thepage}
}

% Header für das Tabellenverzeichnis
\fancypagestyle{lot}{
  \fancyhf{}
  \renewcommand\headrulewidth{0.4pt}
  \fancyhead[RE,LO]{LIST OF TABLES}
  \fancyhead[RO,LE]{\thepage}
}

% Header für das Abkürzungsverzeichnis
\fancypagestyle{akv}{
  \fancyhf{}
  \renewcommand\headrulewidth{0.4pt}
  \fancyhead[RE,LO]{LIST OF ABBREVIATIONS}
  \fancyhead[RO,LE]{\thepage}
}

% Header für den Anhang
\fancypagestyle{app}{
  \fancyhf{}
  \renewcommand\headrulewidth{0.4pt}
  \fancyhead[RE,LO]{APPENDIX}
  \fancyhead[RO,LE]{\thepage}
}

% Header für das Abbildungsverzeichnis
\fancypagestyle{abv}{
  \fancyhf{}
  \renewcommand\headrulewidth{0.4pt}
  \fancyhead[RE,LO]{LIST OF FIGURES}
  \fancyhead[RO,LE]{\thepage}
}

% Header für das Literaturverzeichnis
\fancypagestyle{bib}{
  \fancyhf{}
  \renewcommand\headrulewidth{0.4pt}
  \fancyhead[RE,LO]{BIBLIOGRAPHY}
  \fancyhead[RO,LE]{\thepage}
}

% Header für den Index
\fancypagestyle{swv}{
  \fancyhf{}
  \renewcommand\headrulewidth{0.4pt}
  \fancyhead[RE,LO]{INDEX}
  \fancyhead[RO,LE]{\thepage}
}

% Header für Listings
\fancypagestyle{lst}{
  \fancyhf{}
  \renewcommand\headrulewidth{0.4pt}
  \fancyhead[RE,LO]{LISTINGS}
  \fancyhead[RO,LE]{\thepage}
}

% -------------------------------------------
% Mathematik-Einstellungen
% -------------------------------------------

% Erlaubt Theorem in schattierter Box
\usepackage{shadethm}

% AMS-LaTeX Erweiterung
\usepackage{amsmath}
\usepackage{amsthm}
\usepackage{amssymb}

% Paket zur Erzeugung von Matrizen
\usepackage{array}

% Equation Counter
\numberwithin{equation}{section}

% Definitionen
\newtheoremstyle{examplestyle}	% Name of the style to be used
  {10mm}						% Measure of space to leave above the theorem. E.g.: 3pt
  {10mm}						% Measure of space to leave below the theorem. E.g.: 3pt
  {}							% Name of font to use in the body of the theorem
  {}							% Measure of space to indent
  {\bfseries}					% Name of head font
  {\newline}					% Punctuation between head and body
  {10mm}						% Space after theorem head
  {}							% Manually specify head

\theoremstyle{examplestyle}

\colorlet{shadethmcolor}{BoxBackground}				% Farbe des Hintergrundes
%\definecolor{shaderulecolor}{rgb}{0.0,0.0,1.0}		% Farbe des Rahmens
%\setlength{\shadeboxrule}{1pt}						% Breite des Rahmens

\newshadetheorem{lemma}{Lemma}[section]
\newshadetheorem{requirement}{Requirement}[section]
\newshadetheorem{example}{Example}[section]

\usepackage{thmtools}

\declaretheoremstyle[
  spaceabove=3pt, 
  spacebelow=3pt, 
  headfont=\bfseries, 
  bodyfont=\normalfont,
  postheadspace=3pt,
  headpunct={\newline}]{definitionstyle}

\declaretheorem[
  mdframed={linewidth=0mm, backgroundcolor=BoxBackground}, 
  style=definitionstyle,
  name=Definition,parent=chapter,
  Refname={Definition,Definitions}]{definition}

% ------------------------------------------
% Code Eingabe
% ------------------------------------------

% Paket für Code Eingabe
\usepackage{listings}

% Überschrift - Listing X.Y: ...
\renewcommand{\lstlistingname}{Algorithmus}

% Standard: Java
\lstset{frame=none,
  language=Java,
  aboveskip=0pt,
  belowskip=0pt,
  showstringspaces=false,
  columns=flexible,
  basicstyle={\small\ttfamily},
  numbers=left,
  numberstyle=\tiny\color{CodeNumbers},
  keywordstyle=\color{CodeKeywords}\bfseries,
  commentstyle=\color{CodeComments}\itshape,
  stringstyle=\color{CodeStrings},
  extendedchars=true,
  backgroundcolor=\color{CodeBackground},
  breaklines=true,
  breakatwhitespace=true,
  tabsize=3
}

% Pseudocode
\lstdefinelanguage{pseudo}{
	keywords={
		in, if, endif, else, for, while, endfor, function, 
		return, break, true, false, when, each, new, boolean, 
		int, List, Set, Stack}, 
	sensitive=true,
	comment=[l]{//}
}

% -----------------------------------------
% Hack um doublepages zu löschen
% -----------------------------------------

\makeatletter
\def\cleardoublepage{\clearpage\thispagestyle{empty}
  \if@twoside \ifodd\c@page\else
  \hbox{}\newpage
  \if@twocolumn\hbox{}\newpage\fi\fi\fi}
\makeatother

% -----------------------------------------
% Abkürzungsverzeichnis
% ----------------------------------------

\usepackage{nomencl}
\let\abk\nomenclature							% Befehl umbenennen in abk
\setlength{\nomlabelwidth}{.20\hsize}			% Punkte zw. Abkürzung und Erklärung
\renewcommand{\nomlabel}[1]{#1 \dotfill}		%
\setlength{\nomitemsep}{-\parsep} 				% Zeilenabstände verkleinern
\makenomenclature

% -------------------------------------
% Sachregister
% ------------------------------------

\usepackage{makeidx}
\makeindex

% -------------------------------------
% Benutzerdefinierte Kommandos
% -------------------------------------

% Für selbst definierte Befehle: \newcommand{\zb}{z.\,B.\xspace} 
% Der Befehl \xspace schaut, was auf den Befehl folgt und fügt ggf. ein Leerzeichen ein.
\usepackage{xspace}

% Wird statt siehe bspw. die Abk. s. bevorzugt:
\renewcommand{\seename}{s.}

% Noch zu Erledigen Hinweis
\newcommand{\todo}[1]{\noindent\begin{small}\textsf{\textcolor{blue}{(\textsc{TODO:} #1)}}\end{small}}

% Für Notizen an der Seite
\newcommand{\notiz}[1]{\marginpar{\tiny#1}}
%\setlength{\marginparwidth}{25mm}
% Gedanken nicht mehr anzeigen
%\renewcommand{\notiz}[1]{}
